\section*{Bab \ref{ch:b1:integration} --- Pengintegralan pada Sebuah Ruas}

\begin{solution}{exo:b1:integration:1}
$\dfrac{1}{x^2-4} = \dfrac{1/4}{x - 2} - \dfrac{1/4}{x+2}$
(lewat penutupan), sehingga
\[
\int_0^1 \frac{\dd x}{x^2 - 4}
= \frac14\Bigl[\ln\abs{x-2} - \ln\abs{x+2}\Bigr]_0^1
= \frac14\Bigl(\ln\frac{1}{3} - \ln 1\Bigr) = -\frac{\ln 3}{4}.
\]

Lewat parsial ($u' = 1$, $v = \ln t$): $\int_1^{\eu} \ln t\,\dd t =
\bigl[t\ln t\bigr]_1^{\eu} - \int_1^{\eu} \dd t = \eu - (\eu - 1) =
1$.

Lewat parsial ($u' = \sin t$, $v = t$): $\int_0^\pi t\sin t\,\dd t =
\bigl[-t\cos t\bigr]_0^\pi + \int_0^\pi \cos t\,\dd t = \pi + 0 =
\pi$.
\end{solution}

\begin{solution}{exo:b1:integration:2}
Lewat penyulihan $u = t^2 + 1$, $\dd u = 2t\,\dd t$:
\[
\int_0^1 \frac{t\,\dd t}{(t^2+1)^2}
= \frac12 \int_1^2 \frac{\dd u}{u^2}
= \frac12\Bigl[-\frac1u\Bigr]_1^2 = \frac14 .
\]
Dengan $u = \sin t$: $\int_0^{\pi/2} \cos^3 t\,\dd t = \int_0^{\pi/2}
(1 - \sin^2 t)\cos t\,\dd t = \bigl[\sin t - \frac{\sin^3
t}{3}\bigr]_0^{\pi/2} = 1 - \frac13 = \frac23$.
\end{solution}

\begin{solution}{exo:b1:integration:3}
$u_n = \dfrac1n \sum_{k=1}^{n} \dfrac{1}{1 + (k/n)^2}$: yaitu \omterm{thm:b1:integration:riemann}{jumlah Riemann}
bagi $x \mapsto \frac{1}{1+x^2}$ pada $\intcc{0}{1}$, sehingga $u_n \to
\int_0^1 \frac{\dd x}{1+x^2} = \arctan 1 = \dfrac\pi4$.

$\ln v_n = \dfrac1n \sum_{k=1}^{n} \ln\dfrac{n+k}{n} = \dfrac 1n
\sum_{k=1}^n \ln\Bigl(1 + \dfrac kn\Bigr) \to \int_0^1 \ln(1+x)\,\dd
x = \bigl[(1+x)\ln(1+x) - x\bigr]_0^1 = 2\ln 2 - 1$. Jadi $v_n \to
\eu^{2\ln 2 - 1} = \dfrac 4\eu$. (Periksa pengenalannya:
$\frac{(2n)!}{n!\,n^n} = \prod_{k=1}^{n} \frac{n+k}{n}$.)
\end{solution}

\begin{solution}{exo:b1:integration:4}
Limitnya $0$. Misalkan $M = \sup \abs f$ dan $\varepsilon \in
\intoo{0}{1}$. Potonglah di $1 - \varepsilon$:
\[
\Bigl| \int_0^1 x^n f \Bigr|
\leq \int_0^{1 - \varepsilon} x^n \abs f + \int_{1-\varepsilon}^1
x^n \abs f
\leq M\,(1-\varepsilon)^n + M\varepsilon .
\]
Karena $(1 - \varepsilon)^n \to 0$ (\cref{exo:b1:seq:3}), maka limsup
ruas kirinya $\leq M\varepsilon$ untuk setiap $\varepsilon$:
sehingga \omterm{thm:b1:integration:def}{integralnya} menuju $0$.
\end{solution}

\begin{solution}{exo:b1:integration:5}
$Q(\lambda) = \int_a^b (f + \lambda g)^2 = \int f^2 + 2\lambda \int
fg + \lambda^2 \int g^2 \geq 0$ untuk setiap $\lambda$. Jika $\int g^2 =
0$, maka $g = 0$ (menurut kepositifan tegasnya,
\cref{thm:b1:integration:props}~(4)) dan ketaksamaannya menjadi $0 \leq
0$. Kalau tidak, $Q$ merupakan kuadratik yang sejati, yang di mana-mana $\geq 0$: sehingga
diskriminannya $\leq 0$, yakni $\bigl(\int fg\bigr)^2 \leq \int
f^2 \int g^2$.

Kesamaannya berlaku jika dan hanya jika diskriminannya lenyap, jika dan hanya jika $Q(\lambda_0) = 0$ untuk
suatu $\lambda_0$, yakni $\int (f + \lambda_0 g)^2 = 0$, yakni (menurut kepositifan
tegasnya lagi) $f = -\lambda_0 g$: jadi kesamaannya berlaku tepat ketika
$f$ dan $g$ sebanding.
\end{solution}

\begin{solution}{exo:b1:integration:6}
Misalkan $\Phi(a) = \int_a^{a+T} f$. Menurut teorema fundamentalnya
(\cref{thm:b1:integration:ftc}), $\Phi$ \omterm{def:b1:derivative:def}{dapat diturunkan} dengan
$\Phi'(a) = f(a + T) - f(a) = 0$: jadi konstan.

Untuk $x > 0$, tulislah $x = nT + r$, $0 \leq r < T$ (dengan $n = \lfloor x/T
\rfloor$). Menurut Chasles:
\[
\int_0^x f = n \int_0^T f + \int_{nT}^{nT + r} f,
\qquad
\Bigl| \int_{nT}^{nT+r} f \Bigr| \leq T \sup_{\intcc{0}{T}} \abs f
= C .
\]
Lalu $\frac 1x \int_0^x f = \frac{nT}{x}\cdot\frac 1T \int_0^T f +
O\bigl(\frac 1x\bigr)$, dan $\frac{nT}{x} \to 1$: sehingga limitnya
$\frac1T \int_0^T f$.
\end{solution}

\begin{solution}{exo:b1:integration:7}
Misalkan $g(x) = f(x) - x$: yang \omterm{def:b1:continuity:continuous}{kontinu}, dengan
\[
\int_0^1 g = \int_0^1 f - \frac12 = 0 .
\]
Seandainya $g$ tak pernah lenyap, maka teorema nilai antara akan memaksa
tanda yang tetap (karena fungsi \omterm{def:b1:continuity:continuous}{kontinu} pada sebuah \omterm{prop:b1:reals:intervals}{selang} yang mengambil kedua tandanya
pasti lenyap); katakanlah $g > 0$. Maka, menurut kepositifan tegasnya
(\cref{thm:b1:integration:props}~(4) yang diterapkan pada $g > 0$, yang memberikan
$\int g > 0$): yang bertentangan dengan $\int g = 0$. Jadi $g(c) = 0$ untuk
suatu $c$: sehingga $f(c) = c$.
\end{solution}

\begin{solution}{exo:b1:integration:8}
\begin{enumerate}
  \item Lewat parsial dengan $u' = \sin t$, $v = \sin^{n-1} t$:
        \[
        W_n = \bigl[-\cos t\sin^{n-1}t\bigr]_0^{\pi/2}
        + (n-1)\int_0^{\pi/2} \cos^2 t\,\sin^{n-2} t\,\dd t
        = (n-1)(W_{n-2} - W_n),
        \]
        sehingga $nW_n = (n-1)W_{n-2}$. Dari $W_0 = \frac\pi2$, $W_1 = 1$:
        \[
        W_{2p} = \frac{(2p-1)(2p-3)\cdots 1}{(2p)(2p-2)\cdots 2}\,
        \frac{\pi}{2} = \frac{(2p)!}{4^p (p!)^2}\,\frac\pi2,
        \qquad
        W_{2p+1} = \frac{(2p)(2p-2)\cdots 2}{(2p+1)(2p-1)\cdots 3}
        = \frac{4^p (p!)^2}{(2p+1)!} .
        \]
  \item Pada $\intoo{0}{\frac\pi2}$, $0 < \sin t < 1$, sehingga $\sin^{n+1} <
        \sin^n$ dan $(W_n)$ turun (secara tegas) dan positif.
        Lalu mengapitnya dengan rekurensinya:
        \[
        \frac{n}{n+1} = \frac{W_{n+1}}{W_{n-1}}
        \leq \frac{W_{n+1}}{W_n} \leq 1
        \quad\implies\quad \frac{W_{n+1}}{W_n} \to 1 .
        \]
        Adapun invariannya: $a_n = (n+1)W_{n+1}W_n$ memenuhi $a_n = a_{n-1}$
        menurut rekurensi $(n+1)W_{n+1} = nW_{n-1}$, sehingga $a_n = a_0 =
        1 \cdot W_1 W_0 = \frac\pi2$. Lalu
        \[
        n W_n^2 \sim (n+1) W_{n+1} W_n = \frac\pi2
        \quad\implies\quad
        W_n \sim \sqrt{\frac{\pi}{2n}} .
        \]
\end{enumerate}
\end{solution}

\begin{solution}{exo:b1:integration:9}
\begin{enumerate}
  \item Pada $\intoo{0}{\pi}$: $x > 0$, $a - bx = b(\frac ab - x) =
        b(\pi - x) > 0$ dan $\sin x > 0$, sehingga integrannya $> 0$
        dan $I_n > 0$ (menurut kepositifan tegasnya). Adapun batasnya: pada
        $\intcc{0}{\pi}$, $x \leq \pi$ dan $a - bx \leq a$, sehingga $P
        \leq \frac{\pi^n a^n}{n!}$ dan $I_n \leq
        \pi\,\frac{(\pi a)^n}{n!}$, yang menuju $0$ (karena
        faktorialnya mengalahkan suku geometrinya: sebab ia suku umum
        deret eksponensial yang konvergen, bandingkan
        \cref{ex:b1:seq:e}); khususnya $I_n < 1$ untuk $n$ yang besar.
  \item Jabarkanlah $x^n(a - bx)^n = \sum_{j=0}^{n} \binom nj a^{n-j}
        (-b)^j x^{n+j}$: sehingga $P = \frac{1}{n!}\sum_j c_j x^{n+j}$
        dengan $c_j$ yang bulat. Lalu $P^{(k)}(0) = 0$ untuk $k < n$
        (menurut valuasinya) dan, untuk $n \leq k \leq 2n$, $P^{(k)}(0) =
        \frac{k!}{n!} c_{k-n}$, yaitu bilangan bulat karena $n! \mid k!$.
        Lebih lanjut $P(\pi - x) = P(x)$ (lewat penyulihan: karena $\pi - x$ menukar
        faktornya, dengan memakai $a - b(\pi - x) = bx$), sehingga $P^{(k)}(\pi)
        = \pm P^{(k)}(0)$: yang bulat pula.
  \item Dengan $Q = P - P'' + P^{(4)} - \dots$ (yang hingga: karena $P$ berderajat
        $2n$): $Q + Q'' = P$, dan
        \[
        \bigl(Q'\sin x - Q\cos x\bigr)' = (Q + Q'')\sin x = P\sin x .
        \]
        Sehingga $I_n = \bigl[Q'\sin x - Q\cos x\bigr]_0^{\pi} = Q(\pi)
        + Q(0)$, yaitu jumlah nilai $P^{(2k)}$ di $0$ dan $\pi$: jadi
        bilangan bulat menurut (2).
  \item Untuk $n$ yang besar, $I_n$ merupakan bilangan bulat dengan $0 < I_n < 1$:
        yang mustahil. Jadi anggapan $\pi = \frac ab$ gagal: sehingga $\pi$
        irasional.
\end{enumerate}
\end{solution}

\begin{solution}{exo:b1:integration:10}
Kasus $C^1$-nya: lewat parsial,
\[
\int_a^b f(t)\sin\lambda t\,\dd t
= \Bigl[-f(t)\frac{\cos\lambda t}{\lambda}\Bigr]_a^b
+ \frac{1}{\lambda}\int_a^b f'(t)\cos\lambda t\,\dd t,
\]
yang nilai mutlaknya terbatas oleh $\frac{2\sup\abs f + (b - a)\sup\abs{f'}}
{\lambda} \to 0$.

Kasus \omterm{def:b1:continuity:continuous}{kontinunya}: misalkan $\varepsilon > 0$ lalu pilihlah \omterm{def:b1:integration:step}{fungsi tangga}
$\varphi$ dengan $\abs{f - \varphi} \leq \varepsilon$
(karena \cref{thm:b1:integration:approx} menyediakan $\varphi \leq f \leq
\psi$ dengan celah $\leq\varepsilon$; ambillah $\varphi$). Maka
\[
\Bigl|\int_a^b f\sin\lambda t\Bigr|
\leq \int_a^b \abs{f - \varphi}
+ \Bigl|\int_a^b \varphi \sin\lambda t\Bigr|
\leq (b-a)\varepsilon
+ \sum_i \abs{c_i}\,\Bigl|\int_{x_{i-1}}^{x_i} \sin\lambda t\,\dd
t\Bigr| ,
\]
dan masing-masing $\bigl|\int \sin \lambda t\bigr| = \bigl|\frac{\cos\lambda
x_{i-1} - \cos\lambda x_i}{\lambda}\bigr| \leq \frac{2}{\lambda}$:
sehingga suku keduanya menuju $0$. Jadi limsupnya $\leq
(b-a)\varepsilon$ untuk setiap $\varepsilon$: sehingga limitnya $0$.
\end{solution}

\begin{solution}{exo:b1:integration:11}
Misalkan $m = \min f$ dan $M = \max f$, yang tercapai menurut teorema nilai
ekstrem. Karena $g \geq 0$: maka $m\,g \leq fg \leq M\,g$, sehingga menurut
kemonotonannya
\[
m \int_a^b g \;\leq\; \int_a^b fg \;\leq\; M \int_a^b g .
\]
Jika $\int_a^b g = 0$: maka kepositifan tegasnya
(\cref{thm:b1:integration:props}~(4)) memaksa $g \equiv 0$, sehingga kedua
ruasnya lenyap, dan sebarang $c$ berlaku. Kalau tidak, $t = \frac{\int
fg}{\int g}$ terletak di $\intcc{m}{M} = f(\intcc{a}{b})$
(\cref{thm:b1:continuity:evt,thm:b1:continuity:ivt}), sehingga $t =
f(c)$ untuk suatu $c$.

Tandanya penting: pada $\intcc{-1}{1}$ dengan $f(t) = g(t) = t$: $\int
fg = \int_{-1}^1 t^2 = \frac23$, sedangkan $f(c)\int_{-1}^1 t\,\dd t
= 0$ untuk setiap $c$.
\end{solution}

\begin{solution}{exo:b1:integration:12}
Jika $f \equiv 0$ maka klaimnya hampa (karena setiap titiknya sebuah nol). Jadi
anggaplah $f \not\equiv 0$ dan andaikan ia mempunyai paling banyak $n$ nol
yang berbeda di $\intoo{a}{b}$; lalu misalkan $z_1 < \dots < z_m$ ($m \leq n$)
nol yang di situ $f$ \emph{berubah tanda} (yang mungkin saja tak ada). Tetapkanlah
$P(t) = \prod_{i=1}^{m}(t - z_i)$ (dengan hasil kali kosong $= 1$), yang
berderajat $m \leq n$. Pada setiap subselang yang dipotong oleh $z_i$, baik
$f$ maupun $P$ bertanda tetap, dan keduanya berbalik tanda ketika melintasi
suatu $z_i$: sehingga hasil kalinya $fP$ bertanda tetap pada seluruh
$\intoo{a}{b}$. Dan karena \omterm{def:b1:continuity:continuous}{kontinu}, tak identik nol, dan
bertanda tetap, maka ia mempunyai $\bigl|\int_a^b fP\bigr| > 0$ (menurut kepositifan
tegasnya yang diterapkan pada $\abs{fP}$). Padahal $\int f P$ merupakan kombinasi
linear momen $\int f\,t^k$, $k \leq n$, yang semuanya nol:
yang bertentangan. Jadi $f$ mempunyai sekurang-kurangnya $n + 1$ nol berbeda di
$\intoo{a}{b}$.
\end{solution}

\begin{solution}{pb:b1:integration:1}
\textbf{1.} Misalkan $N = \lceil 2c \rceil$. Untuk $n \geq N$:
$\frac{c^{n+1}/(n+1)!}{c^n/n!} = \frac{c}{n+1} \leq \frac12$, sehingga
$0 < \frac{c^n}{n!} \leq \frac{c^N}{N!}\,2^{-(n - N)} \to 0$:
jadi lewat apitan.

\textbf{2.} Tetapkanlah $k$; lalu induksi pada $m$. Untuk $m = 0$: $\int_0^1
x^k = \frac{1}{k+1} = \frac{k!\,0!}{(k+1)!}$. Adapun langkahnya, lewat parsial
($u = (1-x)^m$, $v' = x^k$):
\[
\int_0^1 x^k (1-x)^m \dd x
= \Bigl[\frac{x^{k+1}}{k+1}(1-x)^m\Bigr]_0^1
+ \frac{m}{k+1}\int_0^1 x^{k+1}(1-x)^{m-1}\dd x
= \frac{m}{k+1}\cdot\frac{(k+1)!\,(m-1)!}{(k+m+1)!} ,
\]
yaitu $\frac{k!\,m!}{(k+m+1)!}$.

\textbf{3.} Dengan $k = m = n$: $\int_0^1 (x(1-x))^n =
\frac{(n!)^2}{(2n+1)!} = \frac{1}{(2n+1)\binom{2n}{n}}$. Dan karena
$x(1-x) \leq \frac14$ pada $\intcc{0}{1}$, maka \omterm{thm:b1:integration:def}{integralnya} $\leq
4^{-n}$, sehingga $\binom{2n}{n} \geq \frac{4^n}{2n+1}$ ---
yang selaras dengan $\binom{2n}{n}^{1/n} \to 4$
(\cref{pb:b1:seq:1}).

\textbf{4.} Integrannya \omterm{def:b1:continuity:continuous}{kontinu}, $\geq 0$, dan positif
pada $\intoo{0}{1}$, sehingga tak identik nol: jadi \omterm{thm:b1:integration:def}{integralnya}
$> 0$ (\cref{thm:b1:integration:props}~(4)). Adapun batas atasnya:
$(x(1-x))^n \leq 4^{-n}$ dan $g \leq \sup\abs g$, lalu
kemonotonannya.

\textbf{5.} Di sini $A_0 = \eu - 1$; $A_1 = [x\eu^x]_0^1 - \int_0^1
\eu^x = \eu - (\eu - 1) = 1$. Lewat parsial: $A_n = [x^n \eu^x]_0^1 -
n\int_0^1 x^{n-1}\eu^x = \eu - n A_{n-1}$. Adapun batasnya: integrannya
positif, sehingga $A_n > 0$; dan $\eu^x \leq \eu$ memberikan $A_n \leq
\eu\int_0^1 x^n = \frac{\eu}{n+1}$.

\textbf{6.} Di sini $A_0 = -1 + 1\cdot\eu$. Jika $A_{n-1} = \alpha_{n-1} +
\beta_{n-1}\eu$ dengan entri yang bulat, maka
\[
A_n = \eu - n\alpha_{n-1} - n\beta_{n-1}\eu
= \underbrace{(-n\,\alpha_{n-1})}_{\alpha_n}
+ \underbrace{(1 - n\,\beta_{n-1})}_{\beta_n}\,\eu ,
\]
yang keduanya bulat.

\textbf{7.} Jika $\eu = \frac pq$: maka $q A_n = q\alpha_n + p\beta_n
\in \Z$, dan $0 < qA_n \leq \frac{q\eu}{n+1} < 1$ untuk $n$ yang besar:
yaitu bilangan bulat yang tegas di antara $0$ dan $1$ --- yang mustahil. Jadi $\eu
\notin \Q$. Adapun pada \cref{exo:b1:seq:9} bilangan bulat yang terperangkap adalah
$q!\,\frac pq - q!\,a_q$; sedangkan di sini ia $qA_n$: jadi perangkap yang sama, dengan
bahan bakar \omterm{thm:b1:integration:def}{integral}.

\textbf{8.} Untuk $n = 0$: $R_0 = \int_0^1 \eu^t = \eu - 1$, sehingga $\eu =
1 + R_0$. Lewat parsial ($u = \eu^t$, $v = -\frac{(1-t)^{n+1}}{n+1}$):
\[
R_n = \frac{1}{n!}\Bigl(\Bigl[-\frac{(1-t)^{n+1}}{n+1}
\eu^t\Bigr]_0^1 + \frac{1}{n+1}\int_0^1 (1-t)^{n+1}\eu^t\Bigr)
= \frac{1}{(n+1)!} + R_{n+1} ,
\]
sehingga rumusnya merambat dari $n$ ke $n + 1$. Adapun batasnya: $1 \leq
\eu^t \leq \eu$ pada $\intcc{0}{1}$ dan $\int_0^1 (1-t)^n =
\frac{1}{n+1}$ memberikan $\frac{1}{(n+1)!} \leq R_n \leq
\frac{\eu}{(n+1)!}$.

\textbf{9.} Di sini $\sum_{k=0}^{9} \frac{1}{k!} = \frac{986410}{362880}
= 2.71828152\dots$, dan
\[
\frac{1}{10!} = 2.76\cdot10^{-7} \leq R_9 \leq
\frac{2.75}{10!} = 7.58\cdot10^{-7} ,
\]
sehingga $2.7182818 \leq \eu \leq 2.7182823$: jadi beserta buktinya, $\eu =
2.718282$ sampai enam desimal (dengan nilai sejatinya $2.7182818\dots$).

\textbf{10.} Di sini $f(1 - x) = \frac{(1-x)^n x^n}{n!} = f(x)$. Pada
$\intoo{0}{1}$: $0 < x(1-x) \leq \frac14$, sehingga $0 < f \leq
\frac{1}{4^n\,n!}$.

\textbf{11.} Di sini $x^n(1-x)^n = \sum_{j=0}^{n} (-1)^j\binom
nj\,x^{n+j}$, sehingga $f = \frac{1}{n!}\sum_j c_j\,x^{n+j}$ dengan
$c_j \in \Z$. Jadi $f^{(k)}(0) = 0$ untuk $k < n$ atau $k > 2n$,
dan untuk $n \leq k \leq 2n$: $f^{(k)}(0) = \frac{k!}{n!}\,
c_{k-n}$, yaitu bilangan bulat karena $n! \mid k!$. Lalu kesetangkupannya memberikan
$f^{(k)}(1) = (-1)^k f^{(k)}(0) \in \Z$.

\textbf{12.} Di sini $b^n \pi^{2n-2k} = b^n\bigl(\frac
ab\bigr)^{\,n-k} = a^{\,n-k}\,b^{\,k} \in \Z$, sehingga $G(0) =
\sum_k (-1)^k a^{n-k} b^k f^{(2k)}(0)$ dan demikian pula $G(1)$ bersifat
bulat menurut pertanyaan 11.

\textbf{13.} Pada $\pi^2 G + G''$, suku ke-$k$ dari $\pi^2 G$
membawa $\pi^{2n-2k+2}f^{(2k)}$ sedangkan suku $j = k - 1$ dari
$G''$ membawa $(-1)^{k-1}\pi^{2n-2k+2}f^{(2k)}$: sehingga semuanya
saling meniadakan kecuali $k = 0$ pada jumlah pertamanya dan $j = n$ pada
yang kedua, yakni
\[
G'' + \pi^2 G = b^n\bigl(\pi^{2n+2} f + (-1)^n f^{(2n+2)}\bigr)
= b^n \pi^{2n+2} f = \pi^2 a^n f
\]
(karena $f$ berderajat $2n$, sehingga $f^{(2n+2)} = 0$; dan $b^n\pi^{2n} =
a^n$). Lalu
\[
\bigl(G'\sin\pi x - \pi G\cos\pi x\bigr)'
= (G'' + \pi^2 G)\sin \pi x = \pi^2 a^n f\sin\pi x .
\]

\textbf{14.} Dengan mengintegralkannya atas $\intcc{0}{1}$:
\[
\pi^2 a^n \int_0^1 f\sin(\pi x)\,\dd x
= \bigl[G'\sin\pi x - \pi G\cos\pi x\bigr]_0^1
= \pi\bigl(G(1) + G(0)\bigr) ,
\]
sehingga $\pi a^n \int_0^1 f\sin\pi x = G(0) + G(1) \in \Z$. Pada
$\intoo{0}{1}$, $f > 0$ dan $\sin\pi x > 0$: sehingga ruas kirinya
positif, jadi $G(0) + G(1) \geq 1$; dan $\sin \leq 1$ beserta
pertanyaan 10 membatasinya oleh $\frac{\pi a^n}{4^n\,n!}$.

\textbf{15.} Menurut pertanyaan 1 (dengan $c = \frac a4$),
$\frac{\pi a^n}{4^n n!} \to 0$: sehingga untuk $n$ yang besar ia $< 1$,
yang bertentangan dengan $G(0) + G(1) \geq 1$. Jadi tak ada pecahan $\frac ab$
yang sama dengan $\pi^2$: itulah teorema Legendre. Dan seandainya $\pi$ rasional,
maka $\pi^2$ pun demikian: sehingga $\pi \notin \Q$ --- dan implikasinya
hanya berjalan ke arah ini (karena $\sqrt2$ irasional dengan kuadrat yang
rasional), dan itulah sebabnya $\pi^2 \notin \Q$ sungguh lebih kuat
daripada \cref{exo:b1:integration:9}.

\textbf{16.} Kebulatannya: pertanyaan 11--12 (yaitu \omterm{def:b1:derivative:def}{turunan}
titik ujungnya); kekecilannya: pertanyaan 10 dan 1; jembatannya: pertanyaan
13--14 (yaitu \omterm{thm:b1:integration:parts}{pengintegralan parsial} rangkap yang berteleskop). Tanpa
$\frac{1}{n!}$, data titik ujungnya tetap bulat (bahkan lebih
mudah), tetapi batasnya menjadi $\frac{\pi a^n}{4^n}$, yang
menuju $0$ hanya ketika $a < 4$ --- padahal setiap calonnya mempunyai $a =
b\,\pi^2 > 9$. Jadi faktorialnyalah yang persis mengalahkan
pertumbuhan geometri $a^n$: tanpa faktorial, tak ada teorema.

\textbf{17.} Untuk $a \leq 10$, bilangan bulat $G(0) + G(1)$ bersifat
positif dan paling banyak $\pi\,(10/4)^n/n!$. Pada $n = 7$:
$2.5^7 = 610.35\dots$, sehingga batasnya $\frac{\pi \times
610.35}{5040} \approx 0.38 < 1$ (sedangkan pada $n = 6$ ia masih
$1.07$): jadi pertentangannya mendarat pada tahap ketujuh,
secara eksplisit.

\textbf{18.} Lewat dua kali \omterm{thm:b1:integration:parts}{pengintegralan parsial}:
\[
\int_0^1 x(1-x)\sin\pi x\,\dd x
= \frac1\pi\int_0^1 (1 - 2x)\cos\pi x\,\dd x
= \frac{2}{\pi^2}\int_0^1 \sin\pi x\,\dd x
= \frac{2}{\pi^2}\cdot\frac{2}{\pi} = \frac{4}{\pi^3}
\]
(karena suku kurungnya lenyap: yaitu $x(1-x)$ di $0, 1$, dan $\sin\pi x$
di $0, 1$). Secara simbolis, teleskop $n = 1$-nya (tanpa anggapan
atas $\pi$) berbunyi $\pi^3\int_0^1 f_1\sin\pi x = -(f_1''(0) +
f_1''(1))$ dengan $f_1 = x(1 - x)$, $f_1'' = -2$: sehingga ruas kanannya $4$
--- jadi kedua perhitungannya bersesuaian.

\textbf{19.} Fungsi $\eu^x$ memecahkan $y' = y$ dan $\sin\pi x$ memecahkan
$y'' = -\pi^2 y$: yaitu persamaan linear berkoefisien konstan,
sehingga \omterm{thm:b1:integration:parts}{pengintegralan parsial} mendaurkan intinya kembali ke dirinya dan
menjaga semua data \omterm{def:b1:topology:closure}{perbatasannya} di dalam $\Z + \Z\eu$ (masing-masing \omterm{def:b1:poly:def}{polinomial}
bulat dalam $\pi^2$). Sifat ketertutupan itulah yang diperlukan
mesinnya. Lalu dihaluskan dengan inti yang disesuaikan bagi beberapa titik
sekaligus, mekanisme yang sama menghasilkan teorema Hermite (bahwa $\eu$
transenden, 1873) dan teorema Lindemann (bahwa $\pi$ transenden,
1882) --- yang di luar jilid ini.

\textbf{20.} Lewat pembagian \omterm{def:b1:poly:def}{polinomial} (atau kalikan balik lalu periksa):
\[
x^4(1-x)^4 = (x^6 - 4x^5 + 5x^4 - 4x^2 + 4)(1 + x^2) - 4 .
\]
Lalu mengintegralkan kesamaan yang terpampang itu dibagi $1 + x^2$:
\[
\int_0^1 \frac{x^4(1-x)^4}{1+x^2}\dd x
= \Bigl(\frac17 - \frac46 + 1 - \frac43 + 4\Bigr)
- 4\arctan 1
= \frac{22}{7} - \pi ,
\]
dengan memakai $\frac17 - \frac23 + 1 - \frac43 + 4 = 3 + \frac17$ dan
$\arctan 1 = \frac\pi4$.

\textbf{21.} Integrannya \omterm{def:b1:continuity:continuous}{kontinu} dan positif pada
$\intoo{0}{1}$: sehingga \omterm{thm:b1:integration:def}{integralnya} $> 0$, jadi $\pi < \frac{22}{7}$.
Lebih lanjut $\frac12 \leq \frac{1}{1+x^2} \leq 1$ pada
$\intcc{0}{1}$ dan $\int_0^1 (x(1-x))^4 = \frac{(4!)^2}{9!} =
\frac{1}{630}$ (pertanyaan 3):
\[
\frac{1}{1260} \leq \frac{22}{7} - \pi \leq \frac{1}{630}
\quad\Longrightarrow\quad
3.14126 < \frac{22}{7} - \frac{1}{630} \leq \pi \leq
\frac{22}{7} - \frac{1}{1260} < 3.14207 .
\]

\textbf{22.} Modulo $x^2 + 1$: $x^2 \equiv -1$, sehingga $x^{4m} =
(x^2)^{2m} \equiv 1$ dan $(1 - x)^2 = 1 - 2x + x^2 \equiv -2x$,
jadi $(1-x)^{4m} \equiv (-2x)^{2m} = 4^m (x^2)^m \equiv
(-4)^m$: sehingga sisanya adalah konstanta $(-4)^m$, dan
hasil baginya $Q_m$ berkoefisien bulat (karena pembagian oleh \omterm{def:b1:poly:def}{polinomial}
bulat yang \omterm{def:b1:poly:def}{monik}). Lalu \omterm{def:b1:arith:divides}{membagi} kesamaannya dengan $1 + x^2$ dan
mengintegralkannya:
\[
J_m := \int_0^1 \frac{(x(1-x))^{4m}}{1+x^2}\dd x
= s_m + (-4)^m\,\frac{\pi}{4},
\qquad s_m = \int_0^1 Q_m \in \Q .
\]
Lalu memecahkannya bagi $\pi$: dengan $r_m = (-1)^{m+1}\,4^{\,1-m} s_m \in
\Q$, $\;\abs{\pi - r_m} = 4^{\,1-m} J_m \leq
4^{\,1-m}\cdot4^{-4m} = 4^{\,1-5m}$. Untuk $m = 1$: $s_1 =
\frac{22}{7}$, $r_1 = \frac{22}{7}$, dengan batas $4^{-4} =
\frac{1}{256}$ --- jadi pertanyaan 20--21 lagi.

\textbf{23.} Di sini $\bigl|\pi - \frac{22}{7}\bigr| = \frac{22}{7} -
\pi \approx 1.26\cdot10^{-3}$, yaitu enam belas kali lebih baik daripada
tolok ukur berorde $2$ yaitu $\frac{1}{7^2} \approx 2.0\cdot10^{-2}$
yang dijamin Dirichlet (\cref{pb:b1:derivative:1}, pertanyaan 4);
dan $\bigl|\pi - \frac{355}{113}\bigr| \approx 2.7\cdot10^{-7}$
mengalahkan $\frac{1}{113^2} \approx 7.8\cdot10^{-5}$ dengan faktor
$\approx 300$. Namun tak ada pertentangan dengan apa pun yang terbukti: karena ketaksamaan
Liouville \emph{membatasi dari bawah} galat hampiran hanya bagi
\omterm{pb:b1:logic:1}{bilangan aljabar}, dan tak ada batas semacam itu bagi $\pi$ yang tersedia pada
tingkat ini --- jadi $\pi$ bebas dihampiri dengan luar biasa
baik.

\textbf{24.} Lemanya: andaikan $x = \frac pq$ dan $0 < \abs{a_n +
b_n x} \to 0$. Maka $\abs{a_n + b_n x} = \frac{\abs{q a_n + p
b_n}}{q}$, dengan $q a_n + p b_n$ merupakan bilangan bulat taknol (yang taknol
karena nilai mutlaknya $> 0$): sehingga $\abs{a_n + b_n x}
\geq \frac1q$ untuk setiap $n$, yang bertentangan dengan kekonvergenannya ke
$0$. Adapun kejadiannya: pertanyaan 7 ($x = \eu$, $a_n = \alpha_n$, $b_n =
\beta_n$); \cref{exo:b1:seq:9} ($x = \eu$ lagi, dengan $a_q =
-q!\sum_{k \leq q}\frac{1}{k!}$, $b_q = q!$); dan
\cref{pb:b1:derivative:1}, pertanyaan 1, merupakan bentuk geometrinya.
Sedangkan pada Bagian III dan pada \cref{exo:b1:integration:9} perangkapnya berjalan
\emph{di dalam} pertentangannya: karena menganggap kerasionalannya mengubah
sebuah ungkapan menjadi bilangan bulat, yang lalu diapit analisisnya
ke dalam $\intoo{0}{1}$ --- jadi asas yang sama, yang ditransposisikan.

\textbf{25.} (i) Kebulatannya tinggal pada kalkulus titik ujung
\omterm{def:b1:poly:def}{polinomialnya} (pertanyaan 6, 11--12), kekecilannya pada batas yang
diberikan kemonotonan $\sup$-nya (pertanyaan 4, 10), dan jembatannya pada
\omterm{thm:b1:integration:parts}{pengintegralan parsial} (pertanyaan 8, 13--14) --- ketiganya
merupakan teorema bab ini. (ii) Adapun \omterm{thm:b1:integration:def}{integralnya} memasok apa yang tak dapat
dipasok teorema nilai rata-rata: yaitu kesamaan yang \emph{persis} antara
objek analitiknya dan data aritmetikanya (jadi kesamaan, bukan sekadar
ketaksamaan dengan $c$ yang tak diketahui), dan itulah sebabnya mesinnya
mencapai $\pi^2$ sedangkan \cref{pb:b1:derivative:1} hanya mencapai
eksponen hampirannya. (iii) Yang terpetik: $\eu \notin \Q$,
$\pi^2 \notin \Q$ (sehingga $\pi \notin \Q$), $\eu = 2.718282$
yang bersertifikat, $\frac{22}{7} - \frac{1}{630} \leq \pi \leq
\frac{22}{7} - \frac{1}{1260}$, dan $\binom{2n}{n} \geq
\frac{4^n}{2n+1}$. (iv) Adapun \omterm{def:b1:topology:closure}{perbatasannya}: Hermite dan Lindemann mendorong
mesin yang sama ke ketransendenan; sedangkan Apéry (1979) menjalankan
perangkap bilangan bulatnya pada $\zeta(3)$ --- jadi mesinnya masih memproduksi
matematika abad kedua puluh.

\end{solution}
